\documentclass[10pt,a4paper]{amsart}
\usepackage[margin=19mm]{geometry}
\usepackage{amsmath,amssymb,mathtools}
\usepackage[scr=rsfs,cal=euler]{mathalfa}
\usepackage{xfrac}
\usepackage[unicode,hidelinks,bookmarks=false]{hyperref}
\newcommand{\ie}{i.e.\ }
\newcommand{\cf}{cf.\ }
\newcommand{\resp}{respectively\ }
\newcommand{\Subsection}{\subsection}
\providecommand{\nobreakdash}{\nobreak\hbox{-}}
\setlength{\emergencystretch}{2em}
\multlinegap0pt
\usepackage{amssymb}
\usepackage[shortlabels]{enumitem}
\usepackage{mathtools}
\usepackage{tikz}
\newtheorem{proposition}{Proposition}
\title{Bures--Kuratowski Metrics and Simplicial Complexes for Completely Bounded Maps}
\author{Remus Floricel, Sarah Plosker, Avner Sadikov}
\date{Source: arXiv:2604.04248v1 (CC BY 4.0)}
\begin{document}
\maketitle

\noindent\emph{Source and licence.} Bures--Kuratowski Metrics and Simplicial Complexes for Completely Bounded Maps. Version arXiv:2604.04248v1 (CC BY 4.0), distributed by arXiv under the Creative Commons Attribution 4.0 International licence, http://creativecommons.org/licenses/by/4.0/, as recorded in the arXiv licence metadata for this exact version and shown on the abstract page https://arxiv.org/abs/2604.04248v1. Full licence notice: \texttt{licenses/arxiv-2604.04248-CC-BY-4.0.txt}.\par\smallskip

\noindent\textit{Editorial scope.} Counted unit: the article's proposition prop:lower-KSW-nonCP (source0.tex lines 574--593). The article's complete original proof of it is delivered with it (source0.tex lines 595--653, one \texttt{proof} environment of 59 lines). No mathematics is written, repaired or re-derived: the statement and the proof are the article's own lines, in the article's own notation, and no retained line is substituted or re-flowed.  No further source-local context is needed: every cross-reference of every retained line resolves inside the two delivered spans.  Nothing was omitted from any retained span, so the package carries no omission record.  No retained line contains a citation, so the wrapper carries no bibliography and no bibliography entry.  No retained line contains a figure environment, an image-inclusion command or drawing code, so no image is delivered and the package declares no asset dependency.  Editorial formatting only, in addition: an AMS \texttt{amsart} preamble that restates the article's own declarations and macros used by the retained lines, namely the environments \texttt{proposition} on separate counters, together with the standard packages those lines need.  The retained environments print in this document as Proposition 1 in order of appearance, whereas the article prints the same items with the numbering of its own full document.  The complete native source of the article as distributed in the arXiv e-print for this exact version is bundled unchanged as \texttt{sources/197870/source0.tex}.
\begin{proposition}
\label{prop:lower-KSW-nonCP}
Fix $(K,\tau)\in\mathfrak U$.
Then for all $\phi,\psi\in \mathcal X$,
\[
\|\phi-\psi\|_{cb}
\le
\sqrt2\,\|\tau\|_{cb}\, c_+(\tau)\,\bigl(\|\phi\|_{cb}^{1/2}+\|\psi\|_{cb}^{1/2}\bigr)\,\delta_{\mathrm{reg}}(\phi,\psi).
\]
Consequently, for every $\theta\in \mathcal C$, $\lambda>0$, $p\in[1,\infty]$, $\alpha\in(0,1]$, and all
$\phi,\psi\in\mathcal X\setminus \mathcal C$,
\[
\beta^{BK}_{\theta,\lambda,p,\alpha}(\phi,\psi)
\ge
\lambda\left(
\frac{\|\phi-\psi\|_{cb}}
{\sqrt2\,\|\tau\|_{cb}\, c_+(\tau)\,\bigl(\|\phi\|_{cb}^{1/2}+\|\psi\|_{cb}^{1/2}\bigr)}
\right)^\alpha.
\]
\end{proposition}

\begin{proof}
Fix $\varepsilon>0$. By definition of the Hausdorff distance $\delta_{\mathrm{reg}}^{\tau}$,
for any $X\in\mathscr F_{\tau}(\phi)$ there exists $Y\in\mathscr F_{\tau}(\psi)$ such that
\[
\|X-Y\|\le \delta_{\mathrm{reg}}^{\tau}(\phi,\psi)+\varepsilon.
\]
Choose $W\in\mathscr R_{\tau}(\phi)$ with $X=\Gamma_{\tau}(W)$, and pick $W'\in\mathscr R_{\tau}(\psi)$ with
$Y=\Gamma_{\tau}(W')$ satisfying the bound. Since $\Gamma_{\tau}$ remembers $W$ in its first coordinate,
\[
\|W-W'\|\le \sqrt2\,\|\Gamma_{\tau}(W)-\Gamma_{\tau}(W')\|
\le \sqrt2\bigl(\delta_{\mathrm{reg}}^{\tau}(\phi,\psi)+\varepsilon\bigr).
\]
Now $\phi(a)=W^*\tau(a)W$ and $\psi(a)=W'^*\tau(a)W'$ for all $a\in A$, hence for any $n$ and any $x\in M_n(A)$,
\[
\|\phi^{(n)}(x)-\psi^{(n)}(x)\|
=
\|(W-W')^*\tau^{(n)}(x)W + W'^*\tau^{(n)}(x)(W-W')\|
\le \|\tau\|_{cb}\,\|x\|\,(\|W\|+\|W'\|)\,\|W-W'\|.
\]
Taking suprema over $n$ and $x$ yields
\[
\|\phi-\psi\|_{cb}
\le \|\tau\|_{cb}\,(\|W\|+\|W'\|)\,\|W-W'\|.
\]
By definition of $\mathscr R_{\tau}(\cdot)$, $\|W\|=\|\phi\|_{cb}^{1/2}$ and $\|W'\|=\|\psi\|_{cb}^{1/2}$.
Letting $\varepsilon\downarrow 0$ gives
\[
\|\phi-\psi\|_{cb}
\le
\sqrt2\,\|\tau\|_{cb}\,\bigl(\|\phi\|_{cb}^{1/2}+\|\psi\|_{cb}^{1/2}\bigr)\,\delta_{\mathrm{reg}}^{\tau}(\phi,\psi).
\]
Since
\[
\delta_{\mathrm{reg}}(\phi,\psi)
=\sup_{(L,\sigma)\in\mathfrak U}\frac{1}{c_+(\sigma)}\,\delta_{\mathrm{reg}}^{\sigma}(\phi,\psi)
\ge \frac{1}{c_+(\tau)}\,\delta_{\mathrm{reg}}^{\tau}(\phi,\psi),
\]
it follows that $\delta_{\mathrm{reg}}^{\tau}(\phi,\psi)\le c_+(\tau)\,\delta_{\mathrm{reg}}(\phi,\psi),$
and therefore
\[
\|\phi-\psi\|_{cb}
\le
\sqrt2\,\|\tau\|_{cb}\, c_+(\tau)\,\bigl(\|\phi\|_{cb}^{1/2}+\|\psi\|_{cb}^{1/2}\bigr)\,\delta_{\mathrm{reg}}(\phi,\psi).
\]
This proves the first inequality.

If now $\phi,\psi\notin\mathcal C$, then by definition of the BK metric,
\[
\beta^{BK}_{\theta,\lambda,p,\alpha}(\phi,\psi)
=\bigl\|(0,\lambda\,\delta_{\mathrm{reg}}(\phi,\psi)^\alpha)\bigr\|_{\ell^p(\mathbb R^2)}
=\lambda\,\delta_{\mathrm{reg}}(\phi,\psi)^\alpha.
\]
Hence
\[
\delta_{\mathrm{reg}}(\phi,\psi)=\lambda^{-1/\alpha}\bigl(\beta^{BK}_{\theta,\lambda,p,\alpha}(\phi,\psi)\bigr)^{1/\alpha},
\]
and substituting this into the first inequality yields the stated estimate on
$\mathcal X\setminus\mathcal C$.
\end{proof}

\end{document}
